\IfFileExists{stacks-project.cls}{%
\documentclass{stacks-project}
}{%
\documentclass{amsart}
}

% For dealing with references we use the comment environment
\usepackage{verbatim}
\newenvironment{reference}{\comment}{\endcomment}
%\newenvironment{reference}{}{}
\newenvironment{slogan}{\comment}{\endcomment}
\newenvironment{history}{\comment}{\endcomment}

% For commutative diagrams we use Xy-pic
\usepackage[all]{xy}

% We use 2cell for 2-commutative diagrams.
\xyoption{2cell}
\UseAllTwocells

% We use multicol for the list of chapters between chapters
\usepackage{multicol}

% This is generally recommended for better output
\usepackage{lmodern}
\usepackage[T1]{fontenc}

% For cross-file-references
\usepackage{xr-hyper}

% Package for hypertext links:
\usepackage{hyperref}

% For any local file, say "hello.tex" you want to link to please
% use \externaldocument[hello-]{hello}
\externaldocument[introduction-]{introduction}
\externaldocument[conventions-]{conventions}
\externaldocument[sets-]{sets}
\externaldocument[categories-]{categories}
\externaldocument[topology-]{topology}
\externaldocument[sheaves-]{sheaves}
\externaldocument[sites-]{sites}
\externaldocument[stacks-]{stacks}
\externaldocument[fields-]{fields}
\externaldocument[algebra-]{algebra}
\externaldocument[brauer-]{brauer}
\externaldocument[homology-]{homology}
\externaldocument[derived-]{derived}
\externaldocument[simplicial-]{simplicial}
\externaldocument[more-algebra-]{more-algebra}
\externaldocument[smoothing-]{smoothing}
\externaldocument[modules-]{modules}
\externaldocument[sites-modules-]{sites-modules}
\externaldocument[injectives-]{injectives}
\externaldocument[cohomology-]{cohomology}
\externaldocument[sites-cohomology-]{sites-cohomology}
\externaldocument[dga-]{dga}
\externaldocument[dpa-]{dpa}
\externaldocument[sdga-]{sdga}
\externaldocument[hypercovering-]{hypercovering}
\externaldocument[schemes-]{schemes}
\externaldocument[constructions-]{constructions}
\externaldocument[properties-]{properties}
\externaldocument[morphisms-]{morphisms}
\externaldocument[coherent-]{coherent}
\externaldocument[divisors-]{divisors}
\externaldocument[limits-]{limits}
\externaldocument[varieties-]{varieties}
\externaldocument[topologies-]{topologies}
\externaldocument[descent-]{descent}
\externaldocument[perfect-]{perfect}
\externaldocument[more-morphisms-]{more-morphisms}
\externaldocument[flat-]{flat}
\externaldocument[groupoids-]{groupoids}
\externaldocument[more-groupoids-]{more-groupoids}
\externaldocument[etale-]{etale}
\externaldocument[chow-]{chow}
\externaldocument[intersection-]{intersection}
\externaldocument[pic-]{pic}
\externaldocument[weil-]{weil}
\externaldocument[adequate-]{adequate}
\externaldocument[dualizing-]{dualizing}
\externaldocument[duality-]{duality}
\externaldocument[discriminant-]{discriminant}
\externaldocument[derham-]{derham}
\externaldocument[local-cohomology-]{local-cohomology}
\externaldocument[algebraization-]{algebraization}
\externaldocument[curves-]{curves}
\externaldocument[resolve-]{resolve}
\externaldocument[models-]{models}
\externaldocument[functors-]{functors}
\externaldocument[equiv-]{equiv}
\externaldocument[pione-]{pione}
\externaldocument[etale-cohomology-]{etale-cohomology}
\externaldocument[proetale-]{proetale}
\externaldocument[relative-cycles-]{relative-cycles}
\externaldocument[more-etale-]{more-etale}
\externaldocument[trace-]{trace}
\externaldocument[crystalline-]{crystalline}
\externaldocument[spaces-]{spaces}
\externaldocument[spaces-properties-]{spaces-properties}
\externaldocument[spaces-morphisms-]{spaces-morphisms}
\externaldocument[decent-spaces-]{decent-spaces}
\externaldocument[spaces-cohomology-]{spaces-cohomology}
\externaldocument[spaces-limits-]{spaces-limits}
\externaldocument[spaces-divisors-]{spaces-divisors}
\externaldocument[spaces-over-fields-]{spaces-over-fields}
\externaldocument[spaces-topologies-]{spaces-topologies}
\externaldocument[spaces-descent-]{spaces-descent}
\externaldocument[spaces-perfect-]{spaces-perfect}
\externaldocument[spaces-more-morphisms-]{spaces-more-morphisms}
\externaldocument[spaces-flat-]{spaces-flat}
\externaldocument[spaces-groupoids-]{spaces-groupoids}
\externaldocument[spaces-more-groupoids-]{spaces-more-groupoids}
\externaldocument[bootstrap-]{bootstrap}
\externaldocument[spaces-pushouts-]{spaces-pushouts}
\externaldocument[spaces-chow-]{spaces-chow}
\externaldocument[groupoids-quotients-]{groupoids-quotients}
\externaldocument[spaces-more-cohomology-]{spaces-more-cohomology}
\externaldocument[spaces-simplicial-]{spaces-simplicial}
\externaldocument[spaces-duality-]{spaces-duality}
\externaldocument[formal-spaces-]{formal-spaces}
\externaldocument[restricted-]{restricted}
\externaldocument[spaces-resolve-]{spaces-resolve}
\externaldocument[formal-defos-]{formal-defos}
\externaldocument[defos-]{defos}
\externaldocument[cotangent-]{cotangent}
\externaldocument[examples-defos-]{examples-defos}
\externaldocument[algebraic-]{algebraic}
\externaldocument[examples-stacks-]{examples-stacks}
\externaldocument[stacks-sheaves-]{stacks-sheaves}
\externaldocument[criteria-]{criteria}
\externaldocument[artin-]{artin}
\externaldocument[quot-]{quot}
\externaldocument[stacks-properties-]{stacks-properties}
\externaldocument[stacks-morphisms-]{stacks-morphisms}
\externaldocument[stacks-limits-]{stacks-limits}
\externaldocument[stacks-cohomology-]{stacks-cohomology}
\externaldocument[stacks-perfect-]{stacks-perfect}
\externaldocument[stacks-introduction-]{stacks-introduction}
\externaldocument[stacks-more-morphisms-]{stacks-more-morphisms}
\externaldocument[stacks-geometry-]{stacks-geometry}
\externaldocument[moduli-]{moduli}
\externaldocument[moduli-curves-]{moduli-curves}
\externaldocument[examples-]{examples}
\externaldocument[exercises-]{exercises}
\externaldocument[guide-]{guide}
\externaldocument[desirables-]{desirables}
\externaldocument[coding-]{coding}
\externaldocument[obsolete-]{obsolete}
\externaldocument[fdl-]{fdl}
\externaldocument[index-]{index}

% Theorem environments.
%
\theoremstyle{plain}
\newtheorem{theorem}[subsection]{Theorem}
\newtheorem{proposition}[subsection]{Proposition}
\newtheorem{lemma}[subsection]{Lemma}

\theoremstyle{definition}
\newtheorem{definition}[subsection]{Definition}
\newtheorem{example}[subsection]{Example}
\newtheorem{exercise}[subsection]{Exercise}
\newtheorem{situation}[subsection]{Situation}

\theoremstyle{remark}
\newtheorem{remark}[subsection]{Remark}
\newtheorem{remarks}[subsection]{Remarks}

\numberwithin{equation}{subsection}

% Macros
%
\def\lim{\mathop{\mathrm{lim}}\nolimits}
\def\colim{\mathop{\mathrm{colim}}\nolimits}
\def\Spec{\mathop{\mathrm{Spec}}}
\def\Hom{\mathop{\mathrm{Hom}}\nolimits}
\def\Ext{\mathop{\mathrm{Ext}}\nolimits}
\def\SheafHom{\mathop{\mathcal{H}\!\mathit{om}}\nolimits}
\def\SheafExt{\mathop{\mathcal{E}\!\mathit{xt}}\nolimits}
\def\Sch{\mathit{Sch}}
\def\Mor{\mathop{\mathrm{Mor}}\nolimits}
\def\Ob{\mathop{\mathrm{Ob}}\nolimits}
\def\Sh{\mathop{\mathit{Sh}}\nolimits}
\def\NL{\mathop{N\!L}\nolimits}
\def\CH{\mathop{\mathrm{CH}}\nolimits}
\def\proetale{{pro\text{-}\acute{e}tale}}
\def\etale{{\acute{e}tale}}
\def\QCoh{\mathit{QCoh}}
\def\Ker{\mathop{\mathrm{Ker}}}
\def\Im{\mathop{\mathrm{Im}}}
\def\Coker{\mathop{\mathrm{Coker}}}
\def\Coim{\mathop{\mathrm{Coim}}}

% Boxtimes
%
\DeclareMathSymbol{\boxtimes}{\mathbin}{AMSa}{"02}

%
% Macros for moduli stacks/spaces
%
\def\QCohstack{\mathcal{QC}\!\mathit{oh}}
\def\Cohstack{\mathcal{C}\!\mathit{oh}}
\def\Spacesstack{\mathcal{S}\!\mathit{paces}}
\def\Quotfunctor{\mathrm{Quot}}
\def\Hilbfunctor{\mathrm{Hilb}}
\def\Curvesstack{\mathcal{C}\!\mathit{urves}}
\def\Polarizedstack{\mathcal{P}\!\mathit{olarized}}
\def\Complexesstack{\mathcal{C}\!\mathit{omplexes}}
% \Pic is the operator that assigns to X its picard group, usage \Pic(X)
% \Picardstack_{X/B} denotes the Picard stack of X over B
% \Picardfunctor_{X/B} denotes the Picard functor of X over B
\def\Pic{\mathop{\mathrm{Pic}}\nolimits}
\def\Picardstack{\mathcal{P}\!\mathit{ic}}
\def\Picardfunctor{\mathrm{Pic}}
\def\Deformationcategory{\mathcal{D}\!\mathit{ef}}
\usepackage{cleveref}
\crefname{originallemma}{original lemma}{original lemmas}
\Crefname{originallemma}{Original lemma}{Original lemmas}
\makeatletter
\newlabel{lemma-ext-from-perfect-into-bounded-QCoh}{{09M4}{source}{}{sourcepointer-1}{}}
\newlabel{lemma-ext-from-perfect-into-bounded-QCoh@cref}{{[originallemma][0][]09M4}{[1][0][]source}{}{}{}}
\makeatother
\crefname{originalsection}{original section}{original sections}
\Crefname{originalsection}{Original section}{Original sections}
\makeatletter
\newlabel{cohomology-section-global-RHom}{{0B6A}{source}{}{sourcepointer-2}{}}
\newlabel{cohomology-section-global-RHom@cref}{{[originalsection][0][]0B6A}{[1][0][]source}{}{}{}}
\makeatother
\crefname{originallemma}{original lemma}{original lemmas}
\Crefname{originallemma}{Original lemma}{Original lemmas}
\makeatletter
\newlabel{cohomology-lemma-dual-perfect-complex}{{08DQ}{source}{}{sourcepointer-3}{}}
\newlabel{cohomology-lemma-dual-perfect-complex@cref}{{[originallemma][0][]08DQ}{[1][0][]source}{}{}{}}
\makeatother
\crefname{originalremark}{original remark}{original remarks}
\Crefname{originalremark}{Original remark}{Original remarks}
\makeatletter
\newlabel{remark-classical-generator}{{0GEF}{source}{}{sourcepointer-4}{}}
\newlabel{remark-classical-generator@cref}{{[originalremark][0][]0GEF}{[1][0][]source}{}{}{}}
\makeatother
\crefname{originalremark}{original remark}{original remarks}
\Crefname{originalremark}{Original remark}{Original remarks}
\makeatletter
\newlabel{derived-remark-check-on-generator}{{0ATH}{source}{}{sourcepointer-5}{}}
\newlabel{derived-remark-check-on-generator@cref}{{[originalremark][0][]0ATH}{[1][0][]source}{}{}{}}
\makeatother
\crefname{originallemma}{original lemma}{original lemmas}
\Crefname{originallemma}{Original lemma}{Original lemmas}
\makeatletter
\newlabel{coherent-lemma-induction-principle}{{08DR}{source}{}{sourcepointer-6}{}}
\newlabel{coherent-lemma-induction-principle@cref}{{[originallemma][0][]08DR}{[1][0][]source}{}{}{}}
\makeatother
\crefname{originallemma}{original lemma}{original lemmas}
\Crefname{originallemma}{Original lemma}{Original lemmas}
\makeatletter
\newlabel{lemma-affine-compare-bounded}{{06Z0}{source}{}{sourcepointer-7}{}}
\newlabel{lemma-affine-compare-bounded@cref}{{[originallemma][0][]06Z0}{[1][0][]source}{}{}{}}
\makeatother
\crefname{originallemma}{original lemma}{original lemmas}
\Crefname{originallemma}{Original lemma}{Original lemmas}
\makeatletter
\newlabel{lemma-direct-summand-of-a-restriction}{{09IQ}{source}{}{sourcepointer-8}{}}
\newlabel{lemma-direct-summand-of-a-restriction@cref}{{[originallemma][0][]09IQ}{[1][0][]source}{}{}{}}
\makeatother
\crefname{originaltheorem}{original theorem}{original theorems}
\Crefname{originaltheorem}{Original theorem}{Original theorems}
\makeatletter
\newlabel{theorem-bondal-van-den-Bergh}{{09IS}{source}{}{sourcepointer-9}{}}
\newlabel{theorem-bondal-van-den-Bergh@cref}{{[originaltheorem][0][]09IS}{[1][0][]source}{}{}{}}
\makeatother
\crefname{originallemma}{original lemma}{original lemmas}
\Crefname{originallemma}{Original lemma}{Original lemmas}
\makeatletter
\newlabel{cohomology-lemma-pushforward-restriction}{{08DF}{source}{}{sourcepointer-10}{}}
\newlabel{cohomology-lemma-pushforward-restriction@cref}{{[originallemma][0][]08DF}{[1][0][]source}{}{}{}}
\makeatother
\crefname{originallemma}{original lemma}{original lemmas}
\Crefname{originallemma}{Original lemma}{Original lemmas}
\makeatletter
\newlabel{cohomology-lemma-pushforward-perfect}{{08DP}{source}{}{sourcepointer-11}{}}
\newlabel{cohomology-lemma-pushforward-perfect@cref}{{[originallemma][0][]08DP}{[1][0][]source}{}{}{}}
\makeatother
\crefname{originallemma}{original lemma}{original lemmas}
\Crefname{originallemma}{Original lemma}{Original lemmas}
\makeatletter
\newlabel{lemma-orthogonal-koszul-first-variant}{{0GEJ}{source}{}{sourcepointer-12}{}}
\newlabel{lemma-orthogonal-koszul-first-variant@cref}{{[originallemma][0][]0GEJ}{[1][0][]source}{}{}{}}
\makeatother
\crefname{originallemma}{original lemma}{original lemmas}
\Crefname{originallemma}{Original lemma}{Original lemmas}
\makeatletter
\newlabel{lemma-bounded-truncation}{{0GEL}{source}{}{sourcepointer-13}{}}
\newlabel{lemma-bounded-truncation@cref}{{[originallemma][0][]0GEL}{[1][0][]source}{}{}{}}
\makeatother
\crefname{originallemma}{original lemma}{original lemmas}
\Crefname{originallemma}{Original lemma}{Original lemmas}
\makeatletter
\newlabel{lemma-quasi-coherence-direct-image}{{08D5}{source}{}{sourcepointer-14}{}}
\newlabel{lemma-quasi-coherence-direct-image@cref}{{[originallemma][0][]08D5}{[1][0][]source}{}{}{}}
\makeatother
\title{Bounded above complexes are detected by a single perfect generator}
\author{The Stacks Project Authors}
\date{stacks-project:perfect.tex:local-snapshot-20261012; GNU FDL 1.2 or later}
\begin{document}
\maketitle
\setcounter{subsection}{1}
\noindent\textit{Editorial note.} All delivered mathematical text is the author's own, copied from the pinned Stacks Project chapter file perfect.tex; no statement or proof step was rewritten, reordered or removed. The counted span is the author environment and its own proof environment. References inside the span are handed to the assembler as navigation stubs that carry the canonical tag of the cited tag only; whether each of them is load-bearing is recorded on the stub itself.\par\medskip
\noindent\textit{Editorial note.} One sentence inside the counted proof reads, verbatim: 'Small detail omitted.' (at the end of the paragraph extending the property from the generator to all perfect complexes). The counted span is otherwise the author's complete proof of the proposition; per the project rule the phrase is recorded here verbatim.\par\medskip
\noindent\textit{Allowed standard prerequisites (editorial).}\begin{itemize}
\item A map out of a perfect complex into a bounded above complex of quasi-coherent modules vanishes in high degrees. invoked at proof lines 10741-10744. Gives the implication from (1) to (2).
\item The Hom groups, the Ext groups and the derived Hom are related by the standard identifications, and duality of a perfect complex relates the tensor expression to the derived Hom. used at proof lines 10745-10754. Shows that conditions (2) to (6) are all equivalent.
\item A perfect complex generating the derived category is a classical generator, and a property inherited by direct sums, summands and shifts and satisfying the two out of three property holds for all perfect complexes. invoked at proof lines 10755-10766. Extends the equivalent conditions from the given generator to every perfect object, giving condition (7).
\item The induction principle on quasi-compact opens reduces the converse to the affine case and a union step, and the truncation statement converts the vanishing against the Koszul complex into a truncation isomorphism. invoked at proof lines 10767-10832. Organises the converse and supplies the truncation used in the union step.
\end{itemize}\medskip
\ifdefined\hypertarget\hypertarget{sourcepointer-1}{}\fi
\noindent\textit{Source reference.} Original lemma 09M4: sources/stacks-perfect/pinned-original-source.tex, source lines 3465--3465. Invoked at proof lines 10741-10744; the shift making it applicable to the bounded above case is inside the counted span.\par\medskip
\ifdefined\hypertarget\hypertarget{sourcepointer-2}{}\fi
\noindent\textit{Source reference.} Original section 0B6A: sources/stacks-chow/tags.txt, source lines 1--1. Cited at proof lines 10745-10748 for the equivalence of (2), (3), (4); the dual identification linking them to (5), (6) is inside the counted span.\par\medskip
\ifdefined\hypertarget\hypertarget{sourcepointer-3}{}\fi
\noindent\textit{Source reference.} Original lemma 08DQ: sources/stacks-chow/tags.txt, source lines 1--1. Used at proof lines 10751-10754 for the passage between the two blocks of conditions; the observation about shifts is inside the counted span.\par\medskip
\ifdefined\hypertarget\hypertarget{sourcepointer-4}{}\fi
\noindent\textit{Source reference.} Original remark 0GEF: sources/stacks-perfect/pinned-original-source.tex, source lines 4506--4506. Cited at proof lines 10756-10758 to apply the classical generator criterion; the inheritance verification is inside the counted span.\par\medskip
\ifdefined\hypertarget\hypertarget{sourcepointer-5}{}\fi
\noindent\textit{Source reference.} Original remark 0ATH: sources/stacks-chow/tags.txt, source lines 1--1. Used at proof lines 10761-10765 to extend the property to all perfect complexes; the verification that the property has the required inheritance is inside the counted span.\par\medskip
\ifdefined\hypertarget\hypertarget{sourcepointer-6}{}\fi
\noindent\textit{Source reference.} Original lemma 08DR: sources/stacks-chow/tags.txt, source lines 1--1. Invoked at proof lines 10767-10769 to organise the converse; the affine case and the union step are inside the counted span.\par\medskip
\ifdefined\hypertarget\hypertarget{sourcepointer-7}{}\fi
\noindent\textit{Source reference.} Original lemma 06Z0: sources/stacks-perfect/pinned-original-source.tex, source lines 454--454. Used at proof lines 10772-10775 to conclude boundedness above in the affine case; the choice of test object is inside the counted span.\par\medskip
\ifdefined\hypertarget\hypertarget{sourcepointer-8}{}\fi
\noindent\textit{Source reference.} Original lemma 09IQ: sources/stacks-perfect/pinned-original-source.tex, source lines 4058--4058. Invoked at proof lines 10777-10778 together with the generator theorem to produce Q; the rest of the union step is inside the counted span.\par\medskip
\ifdefined\hypertarget\hypertarget{sourcepointer-9}{}\fi
\noindent\textit{Source reference.} Original theorem 09IS: sources/stacks-perfect/pinned-original-source.tex, source lines 4111--4111. Cited with the previous label at proof lines 10777-10778; the Koszul construction that follows is inside the counted span.\par\medskip
\ifdefined\hypertarget\hypertarget{sourcepointer-10}{}\fi
\noindent\textit{Source reference.} Original lemma 08DF: sources/stacks-chow/tags.txt, source lines 1--1. Cited at proof lines 10791-10794 for the two descriptions of the auxiliary complex; the support computation is inside the counted span.\par\medskip
\ifdefined\hypertarget\hypertarget{sourcepointer-11}{}\fi
\noindent\textit{Source reference.} Original lemma 08DP: sources/stacks-chow/tags.txt, source lines 1--1. Cited with the previous label at proof lines 10793-10795 to see that the auxiliary complex is perfect; the use of its second description is inside the counted span.\par\medskip
\ifdefined\hypertarget\hypertarget{sourcepointer-12}{}\fi
\noindent\textit{Source reference.} Original lemma 0GEJ: sources/stacks-perfect/pinned-original-source.tex, source lines 10584--10584. Applied at proof lines 10802-10806 to obtain the truncation isomorphism with the direct image from the intersection; producing the vanishing it needs is inside the counted span.\par\medskip
\ifdefined\hypertarget\hypertarget{sourcepointer-13}{}\fi
\noindent\textit{Source reference.} Original lemma 0GEL: sources/stacks-perfect/pinned-original-source.tex, source lines 10673--10673. Used twice at proof lines 10807-10812 to transfer the vanishing to the first open; the identification of the truncations is inside the counted span.\par\medskip
\ifdefined\hypertarget\hypertarget{sourcepointer-14}{}\fi
\noindent\textit{Source reference.} Original lemma 08D5: sources/stacks-perfect/pinned-original-source.tex, source lines 661--661. Cited at proof lines 10827-10830 to conclude boundedness above on the whole scheme; the final deduction is inside the counted span.\par\medskip
\begin{proposition}
\label{proposition-detecting-bounded-above}
Let $X$ be a quasi-compact and quasi-separated scheme. Let
$G \in D_{perf}(\mathcal{O}_X)$ be a perfect complex which generates 
$D_\QCoh (\mathcal{O}_X)$. Let $E \in D_\QCoh (\mathcal{O}_X)$.
The following are equivalent
\begin{enumerate}
\item $E \in D^-_\QCoh (\mathcal{O}_X)$,
\item $\Hom_{D(\mathcal{O}_X)}(G[-i], E) = 0$ for $i \gg 0$,
\item $\Ext^i_X(G, E) = 0$ for $i \gg 0$,
\item $R\Hom_X(G, E)$ is in $D^-(\mathbf{Z})$,
\item $H^i(X, G^\vee \otimes_{\mathcal{O}_X}^\mathbf{L} E) = 0$
for $i \gg 0$,
\item $R\Gamma(X, G^\vee \otimes_{\mathcal{O}_X}^\mathbf{L} E)$
is in $D^-(\mathbf{Z})$,
\item for every perfect object $P$ of $D(\mathcal{O}_X)$
\begin{enumerate}
\item the assertions (2), (3), (4) hold with $G$ replaced by $P$, and
\item $H^i(X, P \otimes_{\mathcal{O}_X}^\mathbf{L} E) = 0$ for $i \gg 0$,
\item $R\Gamma(X, P \otimes_{\mathcal{O}_X}^\mathbf{L} E)$
is in $D^-(\mathbf{Z})$.
\end{enumerate}
\end{enumerate}
\end{proposition}

\begin{proof}
Assume (1). Since
$\Hom_{D(\mathcal{O}_X)}(G[-i], E) = \Hom_{D(\mathcal{O}_X)}(G, E[i])$
we see that this is zero for $i \gg 0$ by
Lemma \ref{lemma-ext-from-perfect-into-bounded-QCoh}. This proves
that (1) implies (2).

\medskip\noindent
Parts (2), (3), (4) are equivalent by the discussion in
Cohomology, Section \ref{cohomology-section-global-RHom}.
Part (5) and (6) are equivalent as $H^i(X, -) = H^i(R\Gamma(X, -))$
by definition. The equivalent conditions (2), (3), (4) are
equivalent to the equivalent conditions (5), (6) by
Cohomology, Lemma \ref{cohomology-lemma-dual-perfect-complex}
and the fact that $(G[-i])^\vee = G^\vee[i]$.

\medskip\noindent
It is clear that (7) implies (2). Conversely, 
let us prove that the equivalent conditions (2) -- (6) imply (7).
Recall that $G$ is a  classical generator for $D_{perf}(\mathcal{O}_X)$ by
Remark \ref{remark-classical-generator}.
For $P \in D_{perf}(\mathcal{O}_X)$ let $T(P)$ be the assertion that
$R\Hom_X(P, E)$ is in $D^-(\mathbf{Z})$.
Clearly, $T$ is inherited by direct sums,
satisfies the 2-out-of-three property for distinguished
triangles, is inherited by direct summands, and is preserved by shifts.
Hence by Derived Categories, Remark \ref{derived-remark-check-on-generator}
we see that (4) implies $T$ holds on all of $D_{perf}(\mathcal{O}_X)$.
The same argument works for all other properties, except that for property
(7)(b) and (7)(c) we also use that $P \mapsto P^\vee$ is a self
equivalence of $D_{perf}(\mathcal{O}_X)$. Small detail omitted.

\medskip\noindent
We will prove the equivalent conditions (2) -- (7) imply (1)
using the induction principle of
Cohomology of Schemes, Lemma \ref{coherent-lemma-induction-principle}.

\medskip\noindent
First, we prove (2) -- (7) $\Rightarrow$ (1) if $X$ is affine.
Set $P = \mathcal{O}_X[0]$. From (7) we obtain $H^i (X, E) = 0$ for $i \gg 0$.
Hence (1) follows since $E$ is determined by $R\Gamma (X, E)$,
see Lemma \ref{lemma-affine-compare-bounded}.

\medskip\noindent
Now assume $X = U \cup V$ with $U$ a quasi-compact open of $X$ and
$V$ an affine open, and assume the implication (2) -- (7) $\Rightarrow$ (1)
is known for the schemes $U$, $V$, and $U \cap V$.
Suppose $E \in D_\QCoh(\mathcal{O}_X)$ satisfies (2) -- (7).
By Lemma \ref{lemma-direct-summand-of-a-restriction} and
Theorem \ref{theorem-bondal-van-den-Bergh} there exists a perfect
complex $Q$ on $X$ such that $Q|_U$ generates $D_\QCoh (\mathcal{O}_U)$. 
Let $f_1, \dots , f_r \in \Gamma (V, \mathcal{O}_V)$
be such that $V \setminus U = V(f_1, \dots , f_r)$ as subsets of
$V$. Let $K \in D_{perf}(\mathcal{O}_V)$ be the object
corresponding to the Koszul complex on $f_1, \dots , f_r$.
Let $K' \in D_{perf}(\mathcal{O}_X)$ be
\begin{equation}
\label{equation-detecting-bounded-above}
K' = R (V \to X)_* K = R (V \to X)_! K,
\end{equation}
see Cohomology, Lemmas \ref{cohomology-lemma-pushforward-restriction} and
\ref{cohomology-lemma-pushforward-perfect}. This is a perfect
complex on $X$ supported on the closed set $X \setminus U \subset V$
and isomorphic to $K$ on $V$. By assumption, we know
$R\Hom_{\mathcal{O}_X}(Q, E)$ and
$R\Hom_{\mathcal{O}_X}(K', E)$ are bounded above.

\medskip\noindent
By the second description of $K'$ in (\ref{equation-detecting-bounded-above})
we have
$$
\Hom_{D(\mathcal{O}_V)}(K[-i], E|_V) = \Hom_{D(\mathcal{O}_X)}(K'[-i], E) = 0
$$
for $i \gg 0$. Therefore, we may apply
Lemma \ref{lemma-orthogonal-koszul-first-variant} to $E|_V$ to
obtain an integer $a$ such that
$\tau_{\geq a}(E|_V) = \tau_{\geq a} R (U \cap V \to V)_* (E|_{U \cap V})$.
Then $\tau_{\geq a} E = \tau_{\geq a} R (U \to X)_* (E |_U)$
(check that the canonical map is an isomorphism after restricting to
$U$ and to $V$). Hence using Lemma \ref{lemma-bounded-truncation}
twice we see that
$$
\Hom_{D(\mathcal{O}_U)}(Q|_U [-i], E|_U) =
\Hom_{D(\mathcal{O}_X)}(Q[-i], R (U \to X)_* (E|_U)) = 0
$$
for $i \gg 0$. Since the Proposition holds for $U$ and the generator
$Q|_U$, we have $E|_U \in D^-_\QCoh(\mathcal{O}_U)$. But then since
the functor $R (U \to X)_*$ preserves $D^-_\QCoh$ 
(by Lemma \ref{lemma-quasi-coherence-direct-image}), we get
$\tau_{\geq a}E \in D^-_\QCoh(\mathcal{O}_X)$. Thus 
$E \in D^-_\QCoh (\mathcal{O}_X)$. 
\end{proof}

\end{document}
